\section*{Chapter \ref{ch:b1:structures} --- Algebraic Structures}

\begin{solution}{exo:b1:structures:1}
\emph{Stability:} $x * y = 1 \iff x + y - xy = 1 \iff (1-x)(1-y) = 0$,
impossible for $x, y \neq 1$. Indeed the key identity is
\[
1 - x * y = (1 - x)(1 - y):
\]
the map $\varphi(x) = 1 - x$ sends $(E, *)$ to $(\R^*, \times)$ with
$\varphi(x * y) = \varphi(x)\varphi(y)$ --- a bijective morphism. All
axioms now transport: associativity and commutativity follow from
those of $\times$; the identity is $\varphi^{-1}(1) = 0$ (check: $x *
0 = x$); the inverse of $x$ is $\varphi^{-1}\bigl((1-x)^{-1}\bigr) =
1 - \frac{1}{1-x} = \frac{x}{x - 1}$ (which is $\neq 1$). So $(E, *)$
is an abelian group.
\end{solution}

\begin{solution}{exo:b1:structures:2}
\begin{enumerate}
  \item Yes: product of positives is positive, identity $1$, inverse
        $\frac 1x$, associativity inherited from $\R^*$.
  \item No: not stable ($1 + 1 = 2 \notin \{-1,0,1\}$).
  \item Yes: the standard example.
  \item No: not stable (odd $+$ odd $=$ even), and no identity ($0$ is
        even).
\end{enumerate}
\end{solution}

\begin{solution}{exo:b1:structures:3}
Write $\mathrm{id}$, the transpositions $\tau_{12}, \tau_{13},
\tau_{23}$ (swapping the two named points), and the cycles $c =
(1\,2\,3)$ (i.e.\ $1 \mapsto 2 \mapsto 3 \mapsto 1$) and $c^2 =
(1\,3\,2)$. The table of $\sigma\rho$ (row $\sigma$, column $\rho$,
apply $\rho$ first):
\begin{center}
\renewcommand{\arraystretch}{1.2}
\begin{tabular}{c|cccccc}
$\sigma\backslash\rho$ & $\mathrm{id}$ & $c$ & $c^2$ & $\tau_{12}$ &
$\tau_{13}$ & $\tau_{23}$ \\
\hline
$\mathrm{id}$ & $\mathrm{id}$ & $c$ & $c^2$ & $\tau_{12}$ &
$\tau_{13}$ & $\tau_{23}$ \\
$c$ & $c$ & $c^2$ & $\mathrm{id}$ & $\tau_{13}$ & $\tau_{23}$ &
$\tau_{12}$ \\
$c^2$ & $c^2$ & $\mathrm{id}$ & $c$ & $\tau_{23}$ & $\tau_{12}$ &
$\tau_{13}$ \\
$\tau_{12}$ & $\tau_{12}$ & $\tau_{23}$ & $\tau_{13}$ & $\mathrm{id}$
& $c^2$ & $c$ \\
$\tau_{13}$ & $\tau_{13}$ & $\tau_{12}$ & $\tau_{23}$ & $c$ &
$\mathrm{id}$ & $c^2$ \\
$\tau_{23}$ & $\tau_{23}$ & $\tau_{13}$ & $\tau_{12}$ & $c^2$ & $c$ &
$\mathrm{id}$ \\
\end{tabular}
\end{center}
Non-commuting pair: $\tau_{12}\tau_{13} = c^2$ while
$\tau_{13}\tau_{12} = c$. (To check one entry:
$\tau_{12}\tau_{13}$ sends $1 \xmapsto{\tau_{13}} 3
\xmapsto{\tau_{12}} 3$, $3 \mapsto 1 \mapsto 2$, $2 \mapsto 2 \mapsto
1$: that is $1 \mapsto 3 \mapsto 2 \mapsto 1$, the cycle $c^2 =
(1\,3\,2)$.)
\end{solution}

\begin{solution}{exo:b1:structures:4}
$H$: $1 \in H$; for $z, w \in H$, $\abs{zw^{-1}} = \abs z / \abs w =
1$: the criterion applies. $\R_+^*$: same, $\abs{xy^{-1}}$ replaced by
positivity. Union: $\iu \in H$ and $2 \in \R_+^*$, but $2\iu$ has
modulus $2 \neq 1$ and is not a positive real: $2\iu \notin H \cup
\R_+^*$, so the union is not stable --- not a subgroup (as predicted
by \cref{exo:b1:structures:6}, neither subgroup contains the other).
\end{solution}

\begin{solution}{exo:b1:structures:5}
Morphism: $\eu^{\iu(\theta + \varphi)} =
\eu^{\iu\theta}\eu^{\iu\varphi}$ (\cref{thm:b1:complex:funceq}).
Kernel: $\eu^{\iu\theta} = 1 \iff \theta \in 2\pi\Z$, so $\ker f =
2\pi\Z \neq \{0\}$: not injective. Image: every unit complex number is
$\eu^{\iu\theta}$ for some $\theta$ (polar form), so
$\operatorname{im} f = \mathbb{U}$, the unit circle. The restriction
of $f$ to a half-open interval of length $2\pi$, such as
$\intco{0}{2\pi}$ or $\intoc{-\pi}{\pi}$, is injective (two angles
with the same image differ by a multiple of $2\pi$, and only one
representative of each class fits in the interval); no interval of
greater length works, since it contains two points at distance
$2\pi$.
\end{solution}

\begin{solution}{exo:b1:structures:6}
Intersection: $e \in H \cap K$, and $x, y \in H \cap K$ gives
$xy^{-1}$ in both $H$ and $K$. Union: if $H \subseteq K$ the union is
$K$, a subgroup (and symmetrically). Conversely, suppose neither
inclusion holds: pick $h \in H \setminus K$ and $k \in K \setminus
H$, and suppose $H \cup K$ were a subgroup; then $hk \in H \cup K$.
If $hk \in H$, then $k = h^{-1}(hk) \in H$: contradiction. If $hk \in
K$, then $h = (hk)k^{-1} \in K$: contradiction. So $H \cup K$ is not
a subgroup.
\end{solution}

\begin{solution}{exo:b1:structures:7}
Note first that $x^2 = e$ means $x^{-1} = x$ for every $x$. Then for
$x, y \in G$:
\[
xy = (xy)^{-1} = y^{-1} x^{-1} = yx ,
\]
using \cref{prop:b1:structures:rules}~(2). So $G$ is abelian.
\end{solution}

\begin{solution}{exo:b1:structures:8}
Units of $\Z/18\Z$: classes coprime to $18 = 2 \times 3^2$:
$\overline 1, \overline 5, \overline 7, \overline{11}, \overline{13},
\overline{17}$. Inverse of $\overline 5$: $5 \times 11 = 55 = 3\times
18 + 1$, so $\overline 5^{-1} = \overline{11}$.

$\overline 5 x = \overline 7$: multiply by $\overline{11}$: $x =
\overline{77} = \overline 5$ (since $77 = 4\times 18 + 5$). Unique
solution.

$\overline 6 x = \overline 3$: the equation $6x \equiv 3 \pmod{18}$
means $18 \mid 6x - 3$. But $6x - 3 = 3(2x - 1)$ is odd, while $18$
is even: an even number cannot divide an odd one. No solution.

$\overline 6 x = \overline{12}$: $6x \equiv 12 \pmod{18} \iff x
\equiv 2 \pmod 3$: solutions $x \in \{\overline 2, \overline 5,
\overline 8, \overline{11}, \overline{14}, \overline{17}\}$ --- six
of them.
\end{solution}

\begin{solution}{exo:b1:structures:9}
$\Z[\sqrt 2]$ contains $0$ and $1$, and is stable under subtraction
and product:
\[
(a + b\sqrt 2)(c + d\sqrt 2) = (ac + 2bd) + (ad + bc)\sqrt 2 ,
\]
so it is a subring of $\R$ (commutativity, associativity,
distributivity are inherited). Unit: $(1 + \sqrt 2)(-1 + \sqrt 2) =
2 - 1 = 1$, so $1 + \sqrt 2$ is invertible with inverse $\sqrt 2 - 1
\in \Z[\sqrt 2]$. Its powers $(1 + \sqrt 2)^n$ are strictly
increasing (the base is $> 1$), hence pairwise distinct, and each is
a unit ($\bigl((1+\sqrt2)^n\bigr)^{-1} = (\sqrt 2 - 1)^n$): the group
of units is infinite.
\end{solution}

\begin{solution}{exo:b1:structures:10}
$x + x = (x + x)^2 = x^2 + x^2 + x^2 + x^2 = 4x^2 = 4x$ --- so $2x =
4x$, giving $2x = 0$, i.e.\ $x + x = 0$ (each element is its own
additive inverse). Then
\[
x + y = (x+y)^2 = x^2 + xy + yx + y^2 = x + xy + yx + y ,
\]
so $xy + yx = 0$, i.e.\ $xy = -yx = yx$ (using $-z = z$). Hence $A$ is
commutative.

Example: on $\mathcal{P}(E)$, define $A + B = (A \cup B) \setminus (A
\cap B)$ (symmetric difference) and $A \times B = A \cap B$. One
checks: $(\mathcal{P}(E), +)$ is an abelian group with identity
$\emptyset$ and each set its own inverse; $\cap$ is associative,
commutative, with identity $E$; distributivity $A \cap (B + C) = (A
\cap B) + (A \cap C)$ holds (an element lies in the left side iff it
is in $A$ and in exactly one of $B, C$). And $A \cap A = A$: every
element is idempotent, as required.
\end{solution}

\begin{solution}{exo:b1:structures:11}
Write the hypothesis for $n$ and $n+1$:
\[
(xy)^{n+1} = x^{n+1} y^{n+1}
\quad\text{and}\quad
(xy)^{n+1} = (xy)(xy)^n = xy\,x^n y^n .
\]
Equating: $x^{n+1} y^{n+1} = x\,y\,x^n\,y^n$; cancel $x$ on the left
and $y^n$ on the right: $x^n y = y x^n$. The same computation one
degree higher ($n+1$ and $n+2$) gives $x^{n+1} y = y x^{n+1}$. Then
\[
y\,x^{n+1} = x^{n+1} y = x\,(x^n y) = x\,y\,x^n ,
\]
and cancelling $x^n$ on the right of $y x \cdot x^n = x y \cdot x^n$:
$yx = xy$. So $G$ is abelian.
\end{solution}
